
\chapter{Анализ современных методов восстановления траектории мобильных роботов при кратковременных потерях сигнала ГНСС}
\label{ch:intro}

\section{Классические методы восстановления траектории}

Классические подходы к восстановлению траектории при кратковременной потере ГНСС-сигнала основаны на вероятностных фильтрах и методах оптимизации. Расширенный фильтр Калмана (РФК) \cite{Cahyadi2023} представляет собой вероятностный фильтр для слияния данных БИНС и ГНСС со сглаживанием, включающий прямой проход для фильтрации и обратный проход для коррекции с учетом конечной и начальной точек. Метод обеспечивает интерпретируемость, низкие вычислительные затраты и глобальную оптимизацию с граничными условиями. Данный метод подвержен ошибкам линеаризации в нелинейных сценариях, зависит от модели шума и накапливает ошибки в длительных потерях сигнала.
Бессигматурный фильтр Калмана (БФК) \cite{Pang2024} использует сигма-точки для аппроксимации нелинейных распределений и интеграции конечной точки в обратный проход. По сравнению с РФК, БФК демонстрирует лучшую обработку нелинейностей и устойчивость к шуму БИНС в потерях сигнала, но характеризуется высокой вычислительной сложностью для больших последовательностей и требует точной инициализации.
Сглаживание на основе метода Рауха-Тунга-Штрибеля (RTS) \cite{Fang2020} представляет собой метод, использующий расширенный фильтр Калмана для прямого прохода и обратного сглаживания, интегрируя начальную и конечную позиции для оптимизации траектории на основе данных БИНС. Метод обеспечивает высокую точность за счет глобальной оптимизации с граничными условиями и снижения дрейфа БИНС в длительных интервалах, однако зависит от точной модели системы и обладает вычислительной сложностью для длинных последовательностей.

\section{Нейросетевые методы на основе рекуррентных и трансформерных архитектур}

Современные подходы используют глубокое обучение для моделирования сложных нелинейных зависимостей. Сети на основе LSTM для слияния данных \cite{Kumar2024} применяют рекуррентные сети для моделирования временных зависимостей IMU, обеспечивая захват нелинейных паттернов и подход, основанный на данных, без явных моделей шума. Однако такие методы требуют больших наборов данных, подвержены риску переобучения и характеризуются низкой интерпретируемостью.

В контексте применения рекуррентных сетей именно для компенсации дрейфа инерциальных измерений при потере сигнала ГНСС заслуживает внимания работа c архитектурой из двух независимых однослойных сетей LSTM \cite{Shin2024LSTM}. Первая сеть ориентацию аппарата, вторая — линейную скорость по последовательностям показаний БИНС. Обучение проводилось на данных высокоточной ГНСС c кинематикой в реальном времени. Валидация результатов выполнялась как на реальном мобильном роботе с МЭМС БИНС, так и в симуляции. При длительности потери сигнала ГНСС более 30~с предложенная LSTM-модель снижала СКО позиции более чем на 90~\% по сравнению с классическим решением на базе фильтра Калмана. 



CNN Трансформер \cite{Khatun2022} использует механизм внимания для глобальных зависимостей в оффлайн-режиме для предсказания скорости/позиции с граничными условиями. Метод эффективен в длинных последовательностях благодаря параллельной обработке, но предъявляет высокие вычислительные требования и зависит от качества обучения.
Гибрид CNN-LSTM \cite{Almassri2022} комбинирует сверточные слои для извлечения признаков и LSTM для временного моделирования, обеспечивая улучшенную обработку пространственно-временных данных, снижение ошибок в зашумленных IMU и детектирование шаблонов движения (шаги, махи руками). Недостатками являются сложность архитектуры и необходимость в GPU для обучения.
Искусственная нейронная сеть для слияния данных IMU/UWB с системой захвата движения \cite{Varey2024} представляет собой оффлайн-метод на основе многослойной ANN с backpropagation (Levenberg–Marquardt), интегрирующей измерения IMU и UWB для реконструкции траектории мобильного робота. Начальная и конечная позиции используются как граничные условия для калибровки и оптимизации позиций (X, Y координаты) в экспериментах с круговыми и случайными движениями. Метод обеспечивает высокую точность позиционирования (MAPE 99\%), предотвращает неограниченный дрейф IMU

Выбор нейронных методов восстановления траектории обусловлен не только их превосходной способностью улавливать сложные нелинейные зависимости в данных IMU, недоступные классическим фильтрам Калмана, но и принципиальными аппаратными преимуществами при реализации на современных бортовых платформах. В отличие от традиционных алгоритмов, требующих интенсивной загрузки основного процессора, нейросетевые модели могут выполняться на специализированном нейроускорителе (NPU), который представляет собой полностью отдельное вычислительное ядро. Благодаря этому восстановление траектории при потере сигнала ГНСС не создаёт дополнительной нагрузки на CPU робота и не влияет на работу основных задач управления, стабилизации и планирования движения. Такой подход особенно эффективен на компактных энергоэффективных платформах, таких как NVIDIA Jetson Orin Nano с мощным встроенным NPU и новейший микроконтроллер STM32N6, оснащённый собственным нейроускорителем Neural-ART. В результате нейронные методы обеспечивают оптимальное сочетание высокой точности, низкого энергопотребления и реального времени работы, что критически важно для малых автономных роботов, выполняющих длительные миссии картографирования и инспекции в условиях нестабильного спутникового сигнала.

\subsection{TransformerEncoder}
Transformer Encoder\cite{Vaswani2017, Guyard2025, Zhu2025} представляет собой энкодер, построенный на механизме внимания, который позволяет модели эффективно учитывать глобальные зависимости в последовательностях данных. В отличие от рекуррентных сетей, этот подход не страдает от проблемы затухания градиентов и может параллельно обрабатывать всю последовательность. В реализации используется линейный слой для встраивания входных данных, позиционное кодирование для сохранения порядка временных шагов и несколько слоев энкодера с механизмами само-внимания и прямого распространения. Граничные условия интегрируются через конкатенацию начальной и конечной позиций к входу, что обеспечивает учет контекста без дополнительных проходов в базовом режиме. Такой метод особенно полезен для задач с нелинейными зависимостями в данных БИНС, где глобальный контекст помогает минимизировать накопленные ошибки. Однако квадратичная сложность может ограничивать применение на очень длинных последовательностях.
\subsection{Двунаправленная LSTM}
Двунаправленная LSTM представляет собой рекуррентную сеть, которая обрабатывает последовательности в обоих направлениях, захватывая контекст как из прошлого, так и из будущего. Это позволяет лучше моделировать временные зависимости в данных БИНС, где накопленные ошибки требуют учета всей траектории. В реализации используется линейный слой для встраивания, за которым следует LSTM-блок с двунаправленным блоком. Граничные условия интегрируются через конкатенацию, но для полной эффективности требуется два прохода в двунаправленном режиме. Такой метод устойчив к шуму и подходит для задач с последовательными данными, но может страдать от затухания градиентов на очень длинных интервалах. LSTM балансирует между сложностью и производительностью, делая её хорошей отправной точкой для сравнения с более современными архитектурами. 

Описанные выше подходы наглядно показывают что различные компактные архитектуры нейронных сетей способны улавливать временные зависимость в зашумленных сигналах акселерометра и гироскопа. В то же время для платформ с высокой динамикой, резкими маневрами, сильными возмущениями подходы и архитектуры основанные только на данных (data-driven) оказываются недостаточно устойчивыми. Data-driven архитектуры не учитывают фундаментальные физические ограничения мобильного робота, что помогает обеспечить физическую согласованность, снизить ошибки предсказания при ограниченном объеме данных. 

\section{Физически информированные нейронные сети в задачах восстановления траектории}


В последние годы заметно вырос интерес к физически информированным нейронным сетям (Physics-Informed Neural Networks, PINN) \cite{Sivtsov2025, Bianchi2024}, которые объединяют возможности глубокого обучения с фундаментальными законами физики. В отличие от классических нейросетей, PINN напрямую встраивают дифференциальные уравнения движения в функцию потерь. Это позволяет модели не только аппроксимировать данные, но и соблюдать физические ограничения даже при ограниченном объёме обучающей выборки. 
В робототехнике такой подход особенно востребован для задач динамического моделирования и навигации \cite{Varey2024, Bianchi2024}. Обзоры показывают, что PINN успешно применяются для идентификации параметров роботов, компенсации неопределённостей и построения робастных контроллеров. Особое внимание уделяется беспилотным летательным аппаратам, где низкая инерция и сильные нелинейности делают традиционные методы идентификации малоэффективными.
Применительно к восстановлению траектории при кратковременных потерях сигнала ГНСС физически информированные сети открывают новые возможности. Включение в обучение уравнений Ньютона для поступательного и вращательного движения, а также реальных данных о тяге двигателей позволяет существенно снизить накопление ошибки дрейфа инерциальной системы. Вместо того чтобы полагаться только на статистические закономерности IMU-данных, модель «знает» физические ограничения конкретного робота — массу, моменты инерции, коэффициенты тяги и сопротивления.
Ключевые работы подтверждают эффективность такого направления. В работе \cite{Sivtsov2025} авторы анализируют современные применения PINN в робототехнике и выделяют их как перспективный инструмент для задач с неполными или зашумлёнными измерениями. Так же в работе \cite{Bianchi2024} продемонстрировали использование PINN для оценки динамической модели квадрокоптера, показав, что встроенные физические законы позволяют получать устойчивые результаты даже при значительном уровне шума и ограниченном объёме данных.
Таким образом, интеграция физического знания в нейросетевую архитектуру представляется логичным шагом развития методов восстановления траектории. В настоящей работе этот подход реализован в рамках модели ФИ-НС ПСПС, где данные о скоростях вращения винтов используются напрямую, а уравнения динамики квадрокоптера входят в состав функции потерь. Такой гибридный метод сочетает высокую вычислительную эффективность последовательностных моделей с физической согласованностью результатов, что особенно важно для бортовых систем мобильных роботов.

\section{Гибридные физически информированные модели на основе последовательностей в задачах инерциальной навигации}

В последние годы в задачах инерциальной навигации активно развиваются гибридные архитектуры, сочетающие последовательностные модели с принципами физически информированных нейронных сетей (PINN). Такие подходы позволяют одновременно учитывать сложные временные зависимости в данных IMU и обеспечивать строгое соблюдение фундаментальных законов механики за счёт встраивания остатков дифференциальных уравнений движения в функцию потерь. Авторы отмечают, что именно это сочетание позволяет эффективно решать проблему быстрого накопления ошибок дрейфа при длительных потерях сигнала ГНСС, где традиционные методы и чисто эмпирические нейронные сети часто оказываются недостаточно эффективными.
Примером такого подхода является работа Tan и др.~\cite{Tan2023}, в которой решается задача оценки полного динамического состояния наземного транспортного средства только по данным IMU. Авторы предлагают гибридную схему, где PINN с ограничениями в виде обыкновенных дифференциальных уравнений кинематики и динамики используется для компенсации дрейфа акселерометра и гироскопа, а полученная оценка передаётся в unscented Kalman filter, работающий на многообразиях. Данный подход обеспечивает высокую точность определения трёхмерной позиции, скорости и ориентации даже при длительных сбоях GNSS. В то же время авторы отмечают в качестве недостатка необходимость использования внешнего фильтра, что усложняет полностью энд-to-энд реализацию на борту.
Аналогичный принцип гибридизации применён в работе Zhou и др.~\cite{Zhou2025}, посвящённой предсказанию раскачивания груза на мобильном кране по инерциальным измерениям. В предложенной модели LSTM-подобная последовательностная структура дополнена физическими остатками уравнений движения. Это позволило повысить физическую согласованность результатов и снизить ошибку по сравнению с обычной LSTM. Однако увеличение вычислительной сложности на длинных последовательностях ограничивает возможность применения модели в реальном времени на ресурс-ограниченных устройствах.
Для задач аэрокосмической тематики представляет интерес работа Bianchi и др.~\cite{Bianchi2024}, в которой решается задача идентификации полной динамической модели квадрокоптера, включая трансляционные и вращательные уравнения движения. Авторы используют PINN на основе многослойного персептрона с элементами последовательностной обработки, встраивая уравнения Ньютона–Эйлера непосредственно в функцию потерь. Предложенный подход продемонстрировал преимущество перед расширенным фильтром Калмана по точности при наличии неопределённостей параметров модели. Основным ограничением остаётся квадратичная вычислительная сложность относительно длины обрабатываемой последовательности, что затрудняет применение на бортовых вычислительных системах.
Задача онлайн-оценки состояния квадрокоптера в условиях ветровых возмущений рассмотрена в работе Liu и др.~\cite{Liu2025}. Авторы предложили архитектуру continual PINN с блоками, близкими по структуре к Transformer, и возможностью непрерывного дообучения. Преимуществом подхода является способность модели адаптироваться к изменяющимся внешним условиям без полной переподготовки. В то же время высокие требования к объёму памяти при обработке длинных IMU-последовательностей остаются существенным ограничением.
Лёгкие решения, ориентированные на работу в реальном времени, представлены в работе Bian и др.~\cite{Bian2025}. Авторы решают задачу оценки кинематики суставов человека по данным IMU с помощью компактной физически информированной нейронной сети. Модель сочетает последовательностную обработку с остатками уравнений движения и обеспечивает ошибку менее 5° при возможности функционирования на мобильных устройствах. Ограничением подхода является относительно небольшое число используемых IMU-датчиков, что снижает его обобщаемость на сложные многозвенные механические системы.
Для задач морской и подводной навигации представляет интерес работа An и др.~\cite{An2025}, в которой решается задача идентификации манёвров судна по инерциальным данным. Предложенная гибридная модель на основе последовательностной архитектуры с физическими ограничениями позволяет оптимизировать режим возбуждения системы для повышения точности идентификации параметров. Данный подход обеспечивает существенное снижение ошибок моделирования по сравнению с традиционными методами. Вместе с тем необходимость предварительного расчёта оптимального возбуждения усложняет практическое внедрение.
Дополнительные примеры подтверждают общую тенденцию. В работе Guang и др.~\cite{Guang2021} рассматривается прямое слияние данных IMU и GPS с помощью LSTM; авторы отмечают простоту реализации, но указывают на быстрый рост дрейфа при отсутствии физических ограничений. В исследовании Shurin и Klein~\cite{Shurin2022} предложен гибрид нейросетевой модели и уравнений движения для решения задачи dead reckoning квадрокоптера; подход демонстрирует высокую точность, однако сопровождается ростом вычислительной сложности. Обобщающий анализ представлен в обзоре Sivtsov и др.~\cite{Sivtsov2025}, где авторы констатируют преимущество гибридных моделей по обобщающей способности, но отмечают отсутствие работ, в которых линейные по сложности архитектуры SSM (Mamba) интегрированы с полноценным механизмом PINN. Наконец, в работе Fu и др.~\cite{Fu202X} показано, что гибридные PINN позволяют решать задачи предсказания динамической нагрузки с высокой интерпретируемостью, однако требуют значительного объёма обучающих данных.
Таким образом, существующие гибридные физически информированные модели на основе последовательностных архитектур демонстрируют заметный прогресс в снижении дрейфа IMU и повышении точности траекторного оценивания. В то же время в рецензируемой литературе практически отсутствуют работы, в которых последовательности с селективными пространствами состояний (архитектуры семейства Mamba/SSM) интегрированы с полноценным механизмом физически информированных нейронных сетей для задач инерциальной навигации. Ни одна из проанализированных публикаций не рассматривает комбинацию линейной по сложности Mamba-архитектуры с физическими остатками и параметрической идентификацией динамики. Указанный пробел определяет перспективное направление исследований, способное объединить высокую эффективность обработки длинных последовательностей со строгой физической согласованностью результатов.